\sectiondate{Mémoire sur la ville de Meyrueis}{1881}

« ...C’est une contrée intéressante : la cinquième ville du Diocèse Elle est habitée par une
noblesse et une bourgeoisie nombreuse. On y compte en ce moment douze Chevaliers de St.
Louis et beaucoup d’autres militaires au service du Roi ; beaucoup de bourgeois vivant
noblement, plusieurs avocats ou gradués ; quatre offices de notaire, ses fabriques d’étoffe de
laine de son cru qui partent presque toutes pour l’étranger ne pouvant être servies que par ses
habitants.

Le commerce qui s’y fait des plus beaux mulets du Royaume et des meilleurs moutons la rend
pour ainsi dire, l’entrepôt de la Provence, des provinces voisines et même de l’Espagne.
C’est la mère nourrice des Cévennes par le blé qu’elle leur fournit.
...Il faudrait supprimer toutes les justices bannerettes... que les places de judicature ne
fussent plus vénales, que ce fut le mérite seul qui put y appeler et que le tarif le plus exact
réglât tous les droits ».

\begin{quote}
\itshape
Mémoire sur la ville de Meyrueis, (Annuaire de la Lozère 1881) ainsi que la prise du Malzieu
le 17 Novembre 1573 Pages 3 et ss. \quad Notices rédigées par Fernand André, Archiviste.
\end{quote}
